% Part: proof-theory
% Chapter: cut-elimination
% Section: midsequent

\documentclass[../../../include/open-logic-section]{subfiles}

\begin{document}

\olfileid{pt}{cut}{mid}

\olsection{मध्य-क्रमवर्ती प्रमेय आणि एरब्राँचे प्रमेय}

कट-निर्मूलन प्रमेयाचा एक महत्त्वाचा
परिणाम म्हणजे मध्य-क्रमवर्ती प्रमेय.
त्यात असे म्हटले आहे: सर्व !!{formula}
अग्र-संख्यापक रूपातील असलेल्या
$\Gamma \Sequent \Delta$ ची
!!a{proof}~$\pi$ असेल, तर अशी
!!a{proof}~$\pi'$ असते की तिच्यात
$S \ident \Pi \Sequent
\Lambda$ ही क्रमवर्ती येते
(तिला \emph{मध्य-क्रमवर्ती} म्हणतात)
आणि~$\pi'$ मध्ये $S$ च्या
वरची सर्व अनुमाने विधानीय,
तर $S$ च्या खालची सर्व अनुमाने
संख्यापकांची असतात.
(!!^a{formula}~$!A$ हे
\emph{अग्र-संख्यापक} (किंवा
\emph{अग्र-संख्यापक रूपातील}) असते,
जर ते पुढील रूपात असेल:
\[\mathsf{Q}_1x_1\dots\mathsf{Q}_nx_n\,!B(x_1,\dots,x_n)\]
येथे प्रत्येक~$\mathsf{Q}_i$
हा $\lforall$ किंवा~$\lexists$
असतो.) मध्य-क्रमवर्ती प्रमेयाचा
एक महत्त्वाचा परिणाम म्हणजे
एरब्राँचे प्रमेय. त्याचे व्यापक
रूप थोडे कठीण आहे; साधारणतः
ते असे सांगते: प्रथम-क्रम
तर्कशास्त्रातील प्रत्येक
!!{provable} !!{sentence}~$!A$
साठी असे विधानीय
सर्वतः सत्य सूत्र~$!A'$
असते की त्यावरून $!A$
!!{prove} करता येते.
$!B$ संख्यापकविरहित
असताना $\lexists[x][!B(x)]$
या रूपातील !!{sentence}s साठी
पुढील विधान आहे:

\begin{thm}[एरब्राँचे प्रमेय]
$!B(x)$ संख्यापकविरहित असू द्या
आणि $\Proves
\lexists[x][!B(x)]$ असे समजा.
मग अशी पदे~$t_1$, \dots, $t_n$
असतात की
$\Proves !B(t_1) \lor \dots \lor !B(t_n)$.
\end{thm}

$!B(t_1) \lor \dots \lor !B(t_n)$
या !!{formula} ला
$\lexists[x][!B(x)]$ चा
\emph{एरब्राँचा विकल्पयोग} म्हणतात.
$!B$ संख्यापकविरहित असल्याने
हा विकल्पयोगही संख्यापकविरहित आहे.
म्हणून तो !!{provable} असेल,
तर कटविरहित सिद्धतेने तो
!!{provable} आहे; परिणामी,
केवळ विधानीय अनुमाने
वापरणारी !!a{proof} आहे.
म्हणून तो सर्वतः सत्य आहे.

एरब्राँच्या प्रमेयातील कठीण
भाग म्हणजे प्रत्येक
!!{provable}
!!{formula}~$\lexists[x][!B(x)]$
साठी एरब्राँचा विकल्पयोग
अस्तित्वात आहे हे दाखवणे.
उलट दिशा सोपी आहे:
$\lexists[x][!B(x)]$ चा
एरब्राँचा विकल्पयोग असेल,
तर हे सूत्र सिद्धतायोग्य आहे.
कारण एरब्राँचा विकल्पयोग
$\lexists[x][!B(x)]$
ला निष्पन्न करतो. उदाहरणार्थ,
$n=2$ असताना~$\Log{G3c}$
मधील ही !!a{proof} पाहा:
\[
\Axiom$!B(t_1) \fCenter \lexists[x][!B(x)], !B(t_1), !B(t_2)$
\Axiom$!B(t_2) \fCenter \lexists[x][!B(x)], !B(t_1), !B(t_2)$
\RightLabel{\LeftR{\lor}}
\BinaryInf$!B(t_1) \lor !B(t_2) \fCenter \lexists[x][!B(x)], !B(t_1), !B(t_2)$
\RightLabel{\RightR{\lexists}}
\UnaryInf$!B(t_1) \lor !B(t_2) \fCenter \lexists[x][!B(x)], !B(t_2)$
\RightLabel{\RightR{\lexists}}
\UnaryInf$!B(t_1) \lor !B(t_2) \fCenter \lexists[x][!B(x)]$
\DisplayProof
\]

$\Sequent
\lexists[x][!B(x)]$ ची
!!a{proof}~$\pi$ दिली असता
त्यातून एरब्राँचा विकल्पयोग
मिळवण्यासाठी, कटविरहित
!!{proof}~$\pi$ मधील सर्व
$\RightR{\lexists}$
अनुमाने शेवटी सलग येतात
हा प्रसंग पाहा:
\[
\AxiomC{}
\RightLabel{$\pi_1$}
\Deduce$ \fCenter \lexists[x][!B(x)], !B(t_1), \dots, !B(t_n)$
\RightLabel{\RightR{\lexists}}
\UnaryInf$ \fCenter \lexists[x][!B(x)], !B(t_2), \dots, !B(t_n)$
\RightLabel{$(n-1) \times \RightR{\lexists}$}
\Deduce$ \fCenter \lexists[x][!B(x)]$
\DisplayProof
\]
~$\pi$ मध्ये
$\lexists[x][!B(x)]$
सक्रिय असलेली हीच सर्व
$\RightR{\lexists}$ अनुमाने
असतील, तर~$\pi_1$
मध्ये इतर संख्यापक अनुमाने
नसतात. कारण $!B(x)$
मध्ये अन्य संख्यापक नाहीत
आणि अंतिम क्रमवर्तीच्या
डावीकडे !!{formula}s नाहीत.
म्हणजे $ \fCenter \lexists[x][!B(x)], !B(t_1),
\dots, !B(t_n)$ ही~$\pi$
ची मध्य-क्रमवर्ती आहे.
शिवाय मध्य-क्रमवर्तीच्या
वर संख्यापक अनुमाने
नसल्यामुळे तिच्या वरच्या
सर्व क्रमवर्तींत
$\lexists[x][!B(x)]$
असते, पण ते सक्रिय किंवा
मुख्य नसते. म्हणून
~$\pi_1$ या !!{proof}
मधून ते काढून टाकता येते
आणि $\fCenter !B(t_1), \dots, !B(t_n)$
याची !!a{proof} मिळते.
मग $\RightR{\lor}$ चा
$n-1$ वेळा उपयोग करून
एरब्राँच्या विकल्पयोगाची
!!a{proof},
$\Sequent !B(t_1) \lor \dots \lor !B(t_n)$,
मिळते.

अडचण अशी की
$\Sequent \lexists[x][!B(x)]$
ची कटविरहित !!{proof}
असली तरी तिची सर्व
\RightR{\lexists} अनुमाने
शेवटी एका सलग गटात
येतात असे गृहीत धरता येत
नाही. त्या \RightR{\lexists}
अनुमानांमध्ये एखादे
$!B(t_i)$ मुख्य असलेली
अनुमाने येऊ शकतात.
म्हणून मध्य-क्रमवर्ती
प्रमेयाची गरज आहे.

\begin{thm}[मध्य-क्रमवर्ती प्रमेय]\ollabel{thm:midsequent}
$\Gamma \Sequent \Delta$ ची
~\Log{G3c} मधील
!!{proof}~$\pi$ असेल आणि
$\Gamma$ व $\Delta$ मधील
सर्व !!{formula}s
अग्र-संख्यापक असतील,
तर अशी !!a{proof}~$\pi'$
असते की तिच्यात
$\Pi \Sequent \Lambda$
ही क्रमवर्ती येते;
तिच्या खालची सर्व
अनुमाने संख्यापकांची,
आणि वरची सर्व अनुमाने
विधानीय असतात.
\end{thm}

\begin{proof}
  ~$\pi$ मधील प्रत्येक संख्यापक
  अनुमानाचे \emph{क्रममान} म्हणजे
  त्याच्या खालच्या विधानीय
  अनुमानांची संख्या असे
  परिभाषित करू; आणि
  $\pi$ चे क्रममान~$o(\pi)$
  म्हणजे त्यातील सर्व
  संख्यापक अनुमानांच्या
  क्रममानांची बेरीज.
  $o(\pi) = 0$ असेल, तर
  कोणत्याही संख्यापक
  अनुमानाच्या खाली विधानीय
  अनुमान नसते आणि काम संपते:
  संख्यापक अनुमानच नसल्यास
  अंतिम क्रमवर्ती मध्य-क्रमवर्ती
  घेता येते. अन्यथा सर्व
  संख्यापक अनुमानांना एकच
  आधारक्रमवर्ती असल्याने
  एक सर्वोच्च संख्यापक
  अनुमान असते; त्याची
  आधारक्रमवर्ती मध्य-क्रमवर्ती असते.

  $o(\pi) > 0$ असेल, तर
  क्रममान $>0$ असलेले सर्वात
  खालचे संख्यापक अनुमान
  निवडा. त्याचे क्रममान $>0$
  असल्याने त्याच्या खाली
  विधानीय अनुमान असलेच पाहिजे.
  ते अशा अनुमानांपैकी
  सर्वात खालचे असल्याने
  त्याचा निष्कर्ष दुसऱ्या
  संख्यापक अनुमानाची
  आधारक्रमवर्ती असू शकत
  नाही: तसे असेल, तर
  त्या दुसऱ्या, अधिक
  खालच्या संख्यापक अनुमानाच्या
  खालीही विधानीय अनुमान
  आले असते. कोणते
  संख्यापक अनुमान आणि
  त्याखाली कोणते विधानीय
  अनुमान आहे, यानुसार
  (अनेक!) प्रसंग ठरतात.
  अंतिम क्रमवर्तीमध्ये
  केवळ अग्र-संख्यापक
  सूत्रे असल्यामुळे,
  संख्यापक अनुमानाचे
  मुख्य सूत्र पुढील
  विधानीय अनुमानात
  सक्रिय असू शकत नाही.
  अन्यथा त्या अनुमानाच्या
  मुख्य सूत्रात विधानीय
  संयोजकाच्या व्याप्तीत
  संख्यापक आला असता.
  उप-!!{formula} गुणधर्मामुळे
  ते अंतिम क्रमवर्तीतील
  एखाद्या !!a{formula} चे
  उप-!!{formula} असावे लागले
  असते. म्हणून संख्यापक
  अनुमानाचे मुख्य
  !!{formula} हे खालच्या
  विधानीय अनुमानाच्या
  संदर्भातील !!a{element} आहे.

  पुढील दोन उदाहरणे पाहू:

प्रथम $\RightR{\lforall}$
हे $\LeftR{\lif}$ च्या
उजव्या आधारक्रमवर्तीमध्ये
संपते असे समजा. मग
~$\pi$ ची संबंधित
उप-!!{proof}
ही उप-!!{proof}~$\pi_1$
मध्ये संपते:
\[
\AxiomC{}
\RightLabel{$\pi_2$}
\Deduce$\Gamma' \fCenter \Delta', \lforall[x][!A(x)], !B$
\AxiomC{}
\RightLabel{$\pi_3$}
\Deduce$!C, \Gamma' \fCenter \Delta', !A(c)$
\RightLabel{\RightR{\lforall}}
\UnaryInf$!C, \Gamma' \fCenter \Delta', \lforall[x][!A(x)]$
\RightLabel{\LeftR{\lif}}
\BinaryInf$!B \lif !C, \Gamma' \fCenter \Delta', \lforall[x][!A(x)]$
\DisplayProof
\]
~$\pi$ नियमित आहे असे
मानता येते; म्हणून
आयगेन चर~$c$ हे~$\pi_3$
च्या बाहेर येत नाही.
विशेषतः ते $!B
\lif !C$ मध्ये किंवा
~$\pi_2$ मध्ये येत नाही.
आता सिद्धता~$\pi_1'$ पाहा:
\[
\AxiomC{}
\RightLabel{$\pi_2'$}
\Deduce$\Gamma' \fCenter \Delta', \lforall[x][!A(x)], !A(c), !B$
\AxiomC{}
\RightLabel{$\pi_3$}
\Deduce$!C, \Gamma' \fCenter \Delta', \lforall[x][!A(x)], !A(c)$
\RightLabel{\LeftR{\lif}}
\BinaryInf$!B \lif !C, \Gamma' \fCenter \Delta', \lforall[x][!A(x)], !A(c)$
\RightLabel{\RightR{\lforall}}
\UnaryInf$!B \lif !C, \Gamma' \fCenter \Delta', \lforall[x][!A(x)], \lforall[x][!A(x)]$
\DisplayProof
\]
येथे~$\pi_2'$ ही
~$\pi_2$ च्या उजवीकडे
~$!A(c)$ जोडून दुर्बलीकरणाने
मिळते. (त्यामुळे नवीन
अनुमाने, विशेषतः क्रममानावर
परिणाम करू शकणारी
संख्यापक अनुमाने, जोडली जात
नाहीत. ~$\pi_2$ मध्ये
आयगेन चरे म्हणून
वापरलेले स्थिरांक~$!A(c)$
मध्ये नसल्याने~$\pi_2$
मधील आयगेन चरांच्या
अटीही मोडत नाहीत.)
दुसऱ्या शाखेला आवश्यक
संदर्भ जोडण्यासाठीही
दुर्बलीकरण वापरता येते.
\RightR{\lforall} ची
आयगेन चराची अट अजूनही
पूर्ण आहे, कारण $c$
हे~$!B \lif !C$
मध्ये येत नाही.

आकुंचनाची ग्राह्यता वापरून
$!B \lif !C, \Gamma' \fCenter \Delta',
\lforall[x][!A(x)]$ याची
!!a{proof}~$\pi_1''$
मिळवतो. $\lforall[x][!A(x)]$
साठी $\RightR{\Contraction}$
ची ग्राह्यता सिद्ध करताना
आणि त्या सिद्धतेत वापरलेल्या
$\RightR{\lforall}$ च्या
व्युत्क्रमणीयतेच्या सिद्धतेत
कोणत्याही अनुमानाचे क्रममान
बदलले नाही: जास्तीत जास्त
अनुमाने काढली गेली.
म्हणून~$\pi_1''$ मध्ये
~$\pi_1$ पेक्षा अधिक
संख्यापक अनुमाने नाहीत,
आणि~$\pi_1''$ मधील
प्रत्येक संख्यापक अनुमानाच्या
खाली त्याच्या~$\pi_1$
मधील संबंधित अनुमानाइतकीच
विधानीय अनुमाने आहेत.
अपवाद फक्त अंतिम
\RightR{\lforall} चा:
त्याच्या खाली~$\pi$
मध्ये \LeftR{\lif} होते,
ते आपण~$\pi_1''$ मध्ये
त्याच्या वर हलवले आहे.
~$\pi$ मधील $\pi_1$
च्या जागी $\pi_1''$ ठेवा.
मिळालेल्या !!{proof}
चे क्रममान कमी आहे, कारण
किमान $\RightR{\lforall}$
अनुमानाचे क्रममान (किमान)~$1$
ने कमी झाले. आता
विगमन गृहीतक लावा.

संख्यापक अनुमान
\RightR{\lexists} असल्याचे
प्रसंग आणखी सोपे आहेत:
आकुंचन लागत नाही आणि
आयगेन चराच्या अटीची
काळजीही करावी लागत नाही.
उदाहरणार्थ,
$\RightR{\lexists}$
नंतर $\RightR{\land}$
येतो असे समजा
(यावेळी डाव्या आधारक्रमवर्तीमध्ये).
संबंधित उप-!!{proof}
ही $\pi_1$ आहे:
\[
\AxiomC{}
\RightLabel{$\pi_2$}
\Deduce$\Gamma' \fCenter \Delta', \lexists[x][!A(x)], !A(t), !B$
\RightLabel{\RightR{\lexists}}
\UnaryInf$\Gamma' \fCenter \Delta', \lexists[x][!A(x)], !B$
\AxiomC{}
\RightLabel{$\pi_3$}
\Deduce$\Gamma' \fCenter \Delta', \lexists[x][!A(x)], !C$
\RightLabel{\RightR{\land}}
\BinaryInf$\Gamma' \fCenter \Delta', \lexists[x][!A(x)], !B \land !C$
\DisplayProof
\]
~$\pi$ मधील $\pi_1$
च्या जागी $\pi_1'$ ठेवा:
\[
\AxiomC{}
\RightLabel{$\pi_2$}
\Deduce$\Gamma' \fCenter \Delta', \lexists[x][!A(x)], !A(t), !B$
\AxiomC{}
\RightLabel{$\pi_3'$}
\Deduce$\Gamma' \fCenter \Delta', \lexists[x][!A(x)], !A(t), !C$
\RightLabel{\RightR{\land}}
\BinaryInf$\Gamma' \fCenter \Delta', \lexists[x][!A(x)], !A(t), !B \land !C$
\RightLabel{\RightR{\lexists}}
\UnaryInf$\Gamma' \fCenter \Delta', \lexists[x][!A(x)], !B \land !C$
\DisplayProof
\]
येथे $\pi_3'$ ही
$\pi_3$ च्या उजवीकडे
~$!A(t)$ जोडून
दुर्बलीकरणाने मिळते.
या दुर्बलीकरणात
~$\pi_3$ मधील
प्रत्येक क्रमवर्तीच्या
उजवीकडे केवळ $!A(t)$
जोडले जाते; म्हणून
कोणत्याही संख्यापक
अनुमानाचे क्रममान बदलत
नाही. ~$\pi'$ आणि
~$\pi$ मधील
संख्यापक अनुमानांच्या
क्रममाने सारखीच आहेत;
फक्त~\RightR{\land}
च्या वर हलवलेल्या
\RightR{\lexists}
अनुमानाचे क्रममान~$1$
ने कमी होते. म्हणून
विगमन गृहीतक लागू होते.
\end{proof}

\begin{prob}
\olref{thm:midsequent} च्या
सिद्धतेतील पुढील प्रसंग
वर्णन करा: \LeftR{\lexists}
हे~\RightR{\lif} च्या
आधारक्रमवर्तीमध्ये संपते;
आणि \LeftR{\lforall}
हे~\LeftR{\lor} च्या
डाव्या आधारक्रमवर्तीमध्ये
संपते.
\end{prob}

\begin{proof}[एरब्राँच्या प्रमेयाची सिद्धता]
$\pi$ चा शेवट
$\Sequent \lexists[x][!A(x)]$
मध्ये होतो असे समजा.
\olref{thm:midsequent} नुसार
$\pi$ चे मध्य-क्रमवर्ती~$S$
असलेल्या !!a{proof}~$\pi'$
मध्ये रूपांतर करता येते;
तिचे रूप पुढीलप्रमाणे असले
पाहिजे: \[\Sequent
\lexists[x][!A(x)], !A(t_1), \dots, !A(t_n).\]
$S$ च्या वरची सर्व
अनुमाने विधानीय असल्याने,
$\lexists[x][!A(x)]$
हे $S$ च्या वरच्या
अनुमानात मुख्य असू शकत नाही.
ते अणुरूप नसल्याने
स्वयंसिद्धकातही मुख्य
असू शकत नाही. शिवाय,
$!A(x)$ संख्यापकविरहित
आहे हे लक्षात घेतल्यास
$\lexists[x][!A(x)]$
हे $!A(t_1)$ चे
योग्य उपसूत्र असू शकत नाही;
म्हणून $S$ च्या वर
ते सक्रियही नसते.
त्यामुळे $S$ वर संपणाऱ्या
उप-!!{proof} मधून
$\lexists[x][!A(x)]$
काढल्यावर
$\Sequent !A(t_1), \dots, !A(t_n)$
याची योग्य !!{proof}
मिळते. त्यातून
$\RightR{\lor}$ ची
$n-1$ अनुमाने लावून
$\Sequent !A(t_1) \lor \dots \lor !A(t_n)$
याची !!a{proof} मिळते.
\end{proof}

% \begin{prob}
% Find a Herbrand disjunction of $\lexists[x][\lforall[y][(!A(x) \lif
% !A(y))]]$ by finding a cut-free !!{proof} of $\Sequent
% \lexists[x][\lforall[y][(!A(x) \lif !A(y))]]$ and applying the
% transformation in \olref{thm:midsequent} to find a midsequent (if
% necessary).
% \end{prob}

\end{document}
